\documentclass[11pt]{article}

\usepackage[margin=1in]{geometry}
\usepackage{amsmath,amssymb,amsthm,mathtools}
\usepackage{graphicx}
\usepackage{microtype}
\usepackage[hidelinks]{hyperref}

\newtheorem{theorem}{Theorem}
\newtheorem{lemma}{Lemma}
\newtheorem{remark}{Remark}
\newcommand{\A}{\mathcal A}
\newcommand{\M}{\mathcal M}
\newcommand{\Lcal}{\mathcal L}
\newcommand{\Ocal}{\mathcal O}

\title{A Closure-Stable Banach-Target Converse\\
to the Bornological Operator Identity}
\author{Candidate partial answer to Remark 5.6 of arXiv:math/0506183}
\date{17 August 2026}

\begin{document}
\maketitle

\begin{center}
\fbox{\parbox{0.92\linewidth}{\textbf{Status.} Candidate substantial
partial result, likely valid, subject to human review.  The converse is proved
for Banach targets whose induced ideal bornology is stable under norm closure.
The unrestricted locally convex and non-closure-stable cases remain open.}}
\end{center}

\section{Source question}

Remark 5.6 (PDF p.~18) of Wong,
\emph{The triangle of operators, topologies, bornologies},
\cite{source}, fixes an operator ideal $\A$ on locally
convex spaces and its generated bornology $\M=\mathcal B(\A)$.  It observes
that $\A$-bornologicality of $Y$ implies
\begin{equation}\label{eq:identity}
 \Lcal^b(X,Y)\cap\Ocal(\M)(X,Y)=\Ocal^b(\M)(X,Y)
 \qquad\text{for every locally convex space }X,
\end{equation}
and asks whether the converse holds.

Here $T\in\Lcal^b(X,Y)$ means that some zero-neighborhood has topologically
bounded image; $T\in\Ocal(\M)(X,Y)$ means that $T$ sends every bounded subset
of $X$ to an $\M(Y)$-bounded set; and
$T\in\Ocal^b(\M)(X,Y)$ means that some zero-neighborhood has
$\M(Y)$-bounded image.

\begin{figure}[ht]
  \centering
  \includegraphics[width=\textwidth]{figures/open_question_crop.png}
  \caption{The preceding topological proof and the bornological converse
  question in Remark 5.6 of the source (PDF p.~18).}
\end{figure}

\section{Result}

Call $\M(Y)$ \emph{norm-closure stable} if the norm closure of every
$\M(Y)$-bounded set is again $\M(Y)$-bounded.

\subsection*{Proof intuition}

Test the operator identity on the underlying vector space of $Y$ equipped
with its strongest locally convex topology.  Its bounded sets live in
finite-dimensional subspaces, so the identity map automatically sends all
bounded sets to ideal-bounded sets.  On the other hand, the Banach unit ball
is an absorbing disk and hence a neighborhood for the strongest topology;
the identity is therefore locally bounded.  The assumed equality forces some
absorbing disk to be ideal-bounded.  Closure stability makes its norm closure
ideal-bounded, and Baire category upgrades that closed absorbing disk to a
norm neighborhood.  It then absorbs every norm-bounded set, exactly giving
$\A$-bornologicality.

\begin{theorem}[Banach, closure-stable converse]\label{thm:main}
Let $\A$ be an operator ideal on locally convex spaces,
$\M=\mathcal B(\A)$, and let $Y$ be a Banach space.  Assume:
\begin{enumerate}
\item $\M(Y)$ is norm-closure stable;
\item identity \eqref{eq:identity} holds for every locally convex space $X$.
\end{enumerate}
Then $Y$ is $\A$-bornological.
\end{theorem}

\begin{proof}
Let $Y_{\max}$ denote the vector space underlying $Y$, equipped with the
strongest locally convex topology, and let
\[
 J:Y_{\max}\longrightarrow Y
\]
be the identity map.  Every seminorm on the vector space is continuous for
the strongest topology, so $J$ is continuous.  The closed norm unit ball
$B_Y$ is absolutely convex and absorbing, hence it is a zero-neighborhood in
$Y_{\max}$.  Its image is norm bounded.  Therefore
$J\in\Lcal^b(Y_{\max},Y)$.

We next show that $J\in\Ocal(\M)(Y_{\max},Y)$.  A bounded subset of a vector
space with the strongest locally convex topology is contained in a
finite-dimensional subspace and is bounded there.  Indeed, if a bounded set
contained an infinite linearly independent sequence $(x_n)$, extend it to a
Hamel basis and assign weights tending to infinity to the $x_n$; the
resulting seminorm is continuous for the strongest topology and is unbounded
on the set.  The finite-dimensional assertion then follows from the usual
equivalence of finite-dimensional norms.

Every finite-dimensional bounded subset of $Y$ is $\M(Y)$-bounded: realize
it inside a scalar multiple of the image of a Euclidean unit ball under a
finite-rank map, and recall that every operator ideal contains all
finite-rank operators.  Thus $J$ sends bounded subsets of $Y_{\max}$ to
$\M(Y)$-bounded subsets.

Apply \eqref{eq:identity} with $X=Y_{\max}$.  We obtain
$J\in\Ocal^b(\M)(Y_{\max},Y)$.  Consequently there is an absolutely convex
zero-neighborhood $U$ in $Y_{\max}$ such that $U=J(U)$ is
$\M(Y)$-bounded.  Every zero-neighborhood is absorbing, so
\[
 Y=\bigcup_{m=1}^{\infty}mU
   \subseteq\bigcup_{m=1}^{\infty}m\overline U=Y,
\]
where closure is in the Banach norm.  By norm-closure stability,
$\overline U$ is $\M(Y)$-bounded.  The sets $m\overline U$ are norm closed;
Baire's theorem implies that one of them has nonempty interior.  Since
$\overline U$ is absolutely convex, it is therefore a norm
zero-neighborhood.

It follows that $B_Y\subseteq r\overline U$ for some $r>0$.  Hence $B_Y$ is
$\M(Y)$-bounded.  Every norm-bounded subset of $Y$ is contained in a scalar
multiple of $B_Y$, so the ordinary Banach bornology and the $\A$-bornology
coincide.  This is precisely that $Y$ is $\A$-bornological.
\end{proof}

\section{Examples and scope}

The closure hypothesis is satisfied for several standard ideal bornologies
on Banach targets:
\begin{itemize}
\item finite-dimensional bounded sets (the finite-rank ideal);
\item relatively norm-compact or precompact sets;
\item relatively weakly compact sets.
\end{itemize}
Thus, in each of these cases, the identity in Remark 5.6 characterizes the
corresponding $\A$-bornological Banach spaces.

The proof isolates the obstruction to the unrestricted question.  Without
closure stability, an $\M$-bounded absorbing disk need not absorb the Banach
unit ball uniformly; it may behave like the unit ball of an
everywhere-defined discontinuous stronger norm.  The argument above does not
show that its norm closure remains in $\M(Y)$.  Nor does it cover a general
nonnormable target, because such a space has no topologically bounded
absorbing disk with which to make the identity from the strongest topology
locally bounded.

\paragraph{Upgrade attempts and novelty search.}
Eight focused attempts checked formal dualization, strongest-topology and
mixed-topology constructions, Baire upgrades, common ideals, and possible
non-closed counterexamples.  The attempt record explains why the closure-free
step remains unresolved.  Bounded searches on 17 August 2026 covered the run
indexes, exact title and arXiv id, Remark 5.6, the verbatim open sentence, and
arXiv queries for the displayed operator identity.  No later explicit answer
was found.

\paragraph{Human-review recommendation.}
Verify the characterization of bounded sets in $Y_{\max}$, the use of
finite-rank maps to place them in $\M(Y)$, and norm-closure stability for any
specific ideal to which the theorem is applied.  The main classification
question is whether closure stability can be removed or violated by a genuine
ideal bornology.

\begin{thebibliography}{9}
\bibitem{source}
N.-C. Wong,
\emph{The triangle of operators, topologies, bornologies},
arXiv:math/0506183v1 (2005), Remark 5.6.
\end{thebibliography}

\end{document}
